\section{Research Questions}
\label{sec:research-questions}

An Abstract Graph program specifies which subgraphs and relations become
features. A cycle operator produces cycle features; a neighborhood-pair
operator produces features describing pairs of neighborhoods. Demonstrating
these expected outcomes would only restate the program definition. We evaluate
the language on more substantive grounds: whether it can reproduce established
graph-kernel feature maps with the correct identities and multiplicities,
whether composition provides a practical way to define additional feature
families, and what those programs cost to compute. These aims lead to two
research questions.

\paragraph{RQ1: Representation coverage and composition.}
Which established graph-kernel feature maps can be expressed faithfully as
Abstract Graph programs, which admit only a bounded proxy, and which require
operators not yet present in the language? For faithful constructions, we test
feature identities and multiplicities against independent implementations.
Composed examples then show how mapped objects define higher-order feature
families.

\paragraph{RQ2: Discrimination--complexity trade-off.}
What runtime, memory, and representation size are required by these feature
programs as graph size, structural radius, enumeration bounds, and composition
depth increase? Discrimination profiles describe the equivalence classes
induced by representative programs, but the main comparison concerns faithful
coverage and computational feasibility rather than rediscovering their declared
structural bias.
